\documentclass[11pt,a4paper]{article}
\usepackage[T1]{fontenc}
\usepackage[utf8]{inputenc}
\usepackage{lmodern}
\usepackage{amsmath,amssymb,amsthm,mathtools}
\usepackage[a4paper,margin=25mm,headheight=15pt,footskip=28pt]{geometry}
\usepackage{microtype}
\usepackage{booktabs,tabularx,array,longtable}
\usepackage{enumitem}
\usepackage{xcolor,fancyhdr}
\usepackage{xurl}
\usepackage{hyperref}
\usepackage[nameinlink,noabbrev]{cleveref}
\usepackage{listings}
\usepackage{multicol}

\definecolor{ink}{HTML}{17212B}
\definecolor{deepblue}{HTML}{214E78}
\definecolor{deepgreen}{HTML}{3D6B47}
\definecolor{deepred}{HTML}{8D3A3A}
\definecolor{softblue}{HTML}{EDF4F8}
\definecolor{softgray}{HTML}{F4F6F7}
\definecolor{rulegray}{HTML}{C7CED4}

\hypersetup{
 colorlinks=true,
 linkcolor=deepblue,
 citecolor=deepgreen,
 urlcolor=deepblue,
 pdftitle={Sharp C-r Spectral Disks for the Rvachev--Thue--Morse Transfer Operator},
 pdfauthor={OpenAI; research manuscript prepared for Vladimir Reshetnikov},
 pdfsubject={Exact essential spectra, endpoint-jet ladders, finite resonances, and optimal operator asymptotics},
 pdfkeywords={Rvachev up-function, Fabius function, transfer operator, essential spectrum, C-r regularity, Fredholm theory, Thue--Morse}
}
\urlstyle{tt}

\pagestyle{fancy}
\fancyhf{}
\fancyhead[L]{\small Sharp $C^r$ spectral disks}
\fancyhead[R]{\small ProveIt research extension}
\fancyfoot[C]{\thepage}
\renewcommand{\headrulewidth}{0.4pt}

\setlength{\parindent}{1.1em}
\setlength{\parskip}{0.28em}
\setlength{\emergencystretch}{2.5em}
\allowdisplaybreaks
\setlist[itemize]{itemsep=0.25em,topsep=0.35em,leftmargin=1.5em}
\setlist[enumerate]{itemsep=0.25em,topsep=0.35em,leftmargin=1.7em}

\newtheorem{theorem}{Theorem}[section]
\newtheorem{lemma}[theorem]{Lemma}
\newtheorem{proposition}[theorem]{Proposition}
\newtheorem{corollary}[theorem]{Corollary}
\newtheorem{claim}[theorem]{Claim}
\theoremstyle{definition}
\newtheorem{definition}[theorem]{Definition}
\newtheorem{question}[theorem]{Question}
\newtheorem{example}[theorem]{Example}
\theoremstyle{remark}
\newtheorem{remark}[theorem]{Remark}
\newtheorem*{status}{Claim boundary}
% ed. (2026-09-30): unnumbered environment for ProveIt editorial notes (the theorem counter is unchanged).
\newtheorem*{ednote}{Editorial note (ProveIt, 2026-09-30)}

\newcommand{\C}{\mathbb C}
\newcommand{\R}{\mathbb R}
\newcommand{\N}{\mathbb N}
\newcommand{\Z}{\mathbb Z}
\newcommand{\T}{\mathbb T}
\newcommand{\Dbar}[1]{\overline{D}(0,#1)}
\newcommand{\ess}{\mathrm{ess}}
\newcommand{\Span}{\operatorname{span}}
\newcommand{\ran}{\operatorname{ran}}
\newcommand{\codim}{\operatorname{codim}}
\newcommand{\spec}{\sigma}
\newcommand{\spece}{\sigma_{\mathrm{ess}}}
\newcommand{\norm}[1]{\left\lVert #1\right\rVert}
\newcommand{\abs}[1]{\left|#1\right|}
\newcommand{\dd}{\,\mathrm d}
\newcommand{\e}{\mathrm e}
\newcommand{\ii}{\mathrm i}
\newcommand{\Lop}{\mathcal L}
\newcommand{\Pop}{\mathcal P}
\newcommand{\Uop}{\mathcal U}
\newcommand{\Jop}{\mathcal J}
\newcommand{\Kop}{\mathcal K}
\newcommand{\id}{\mathrm I}
\newcommand{\one}{\mathbf 1}
\newcommand{\Fourier}[1]{\widehat{#1}}
\newcommand{\repo}{\textsc{ProveIt}}
\newcommand{\CrT}{C^r(\T)}
\newcommand{\CrI}{C^r([0,1])}

\newcolumntype{Y}{>{\raggedright\arraybackslash}X}
\newcolumntype{L}[1]{>{\raggedright\arraybackslash}p{#1}}

\lstset{
 basicstyle=\ttfamily\small,
 keywordstyle=\bfseries\color{deepblue},
 commentstyle=\itshape\color{deepgreen},
 frame=single,
 rulecolor=\color{rulegray},
 backgroundcolor=\color{softgray},
 breaklines=true,
 keepspaces=true,
 showstringspaces=false,
 columns=fullflexible,
 xleftmargin=0.5em,xrightmargin=0.5em,
 aboveskip=6pt,belowskip=6pt
}

\title{\textbf{Sharp $C^r$ Spectral Disks for the\\Rvachev--Thue--Morse Transfer Operator}\\[0.55em]
\large Universal dyadic Markov disks, endpoint-jet ladders,\\
finite resonances, and optimal operator asymptotics}
\author{Research manuscript prepared for Vladimir Reshetnikov\\
\small OpenAI}
\date{30 September 2026 (Pacific time)}

\begin{document}
\color{ink}
\maketitle

\begin{abstract}
The Fourier-decay branch of the \repo{} repository asks for the spectrum of
its Rvachev--Thue--Morse transfer operator on finite-smoothness spaces.  The
analytic Hardy-space spectrum had collapsed to three points, and a later
Sobolev manuscript found a regularity-dependent essential disk, but the
periodic H\"older and $C^r$ cases were left open.  This article resolves the
integer $C^r$ problem exactly, both on the circle and on the interval, and
proves a more general structural theorem for smooth dyadic Markov transfer
operators.

For a smooth weight $0\leq\omega\leq1$ satisfying
$\omega(x+1)=1-\omega(x)$, define
\[
 (\Pop_\omega f)(x)=\omega(x)f(x/2)
 +(1-\omega(x))f((x+1)/2).
\]
On $C^r(\T)$, with $r\in\N\cup\{0\}$, its Fredholm essential spectrum is
\emph{always} the exact disk $\Dbar{2^{-r}}$.  Every point of the open disk
is an eigenvalue of infinite geometric multiplicity.  The proof is
constructive: the doubling composition operator is a right inverse, an
explicit infinite-dimensional branch-cancellation space supplies defect
vectors, and differentiation identifies the Calkin image at order $r$ with
$2^{-r}$ times the continuous-space image.  Exterior generalized
eigenvectors automatically bootstrap to $C^\infty$.

On $C^r([0,1])$, unmatched endpoint jets contribute the explicit finite
ladder $\{\omega(0)2^{-j}:0\leq j\leq r\}$ and no other spectral values.
For the Rvachev sine mask $\omega(x)=\sin^2(\pi x/2)$ this ladder collapses
to zero.  Writing $\Lop=\Pop_\omega/2$, we obtain the complete formula
\[
 \spece(\Lop;C^r)=\Dbar{2^{-r-1}},\qquad
 \spec(\Lop;C^r)=\Dbar{2^{-r-1}}\cup\{1/2,-1/4\}
\]
on both the circle and the interval.  The leading resonance is simple; once
$r\ge2$, the resonance $-1/4$ is isolated, semisimple, and two-dimensional.
For $r\ge2$ this gives a sharp two-resonance operator expansion, while a
Calkin-algebra lower bound proves that no compact correction can improve the
residual exponential rate $2^{-r-1}$.  The results are conventional
mathematical proofs in an unrefereed manuscript; no new Lean declaration or
claim of global publication priority is made.
\end{abstract}

\noindent\textbf{Keywords.}
Rvachev up-function; Fabius function; transfer operator; essential spectrum;
Fredholm theory; $C^r$ regularity; endpoint jets; Thue--Morse; spectral
resonance.

\begin{center}
\fcolorbox{deepblue}{softblue}{%
\begin{minipage}{0.91\textwidth}
\begin{status}
The paper proves exact theorems relative to the operator and function spaces
stated here.  It resolves the integer $C^r$ portion of a question explicitly
left open in the inspected \repo{} branch.  It does not claim that every
specialization is globally new, does not settle noninteger H\"older spaces,
and does not constitute a machine-checked formalization.  The repository pin,
source blobs, external analytic inputs, and finite verification scope are
recorded in \cref{app:provenance,app:verification}.
\end{status}
\end{minipage}}
\end{center}

\clearpage
{\small\setlength{\parskip}{0pt}\tableofcontents}
\clearpage

\section{The repository question and the result}
\label{sec:introduction}

\subsection{A delimited open target}
The \repo{} Fabius-function research branch studies the centered Rvachev
up-function, its infinite sinc product, Thue--Morse identities, exact moments,
and several transfer operators.  The manuscript
\emph{Spectral Collapse and Alternating RMS Asymptotics for the Rvachev
Up-Function} proves a finite analytic spectrum and then asks for the complete
spectrum or essential spectral radius on $C^r$, H\"older, Sobolev, and
endpoint-weighted spaces.  Its warning is mathematically important: an
analytic Fredholm determinant cannot simply be transferred to a weaker norm.
A subsequent manuscript answers the periodic Sobolev part and the
integer-order interval Sobolev part, but explicitly leaves H\"older and
$C^r$ spaces untreated \cite{repo-spectral,repo-sobolev}.

This article answers the integer $C^r$ problem.  The answer is not the
analytic three-point set.  At every finite differentiability order there is a
full essential disk.  Its radius is particularly simple in the supremum-jet
geometry of $C^r$: one factor $1/2$ comes from the repository normalization,
and one factor $2^{-r}$ comes from the contraction of the highest derivative.
The two algebraic resonances survive outside that disk exactly when their
moduli exceed this threshold.

The proof also isolates a general mechanism.  It applies to every smooth
nonnegative dyadic partition of unity, not only to the sine mask.  The
essential disk comes from three pieces that fit exactly:
\begin{enumerate}
\item a right inverse generated by the doubling map;
\item an infinite-dimensional branch-cancellation kernel;
\item a derivative quotient that conjugates the Calkin image at order $r$
      to $2^{-r}$ times the continuous-space image.
\end{enumerate}
The sine mask is then special only in the finite resonance calculation and in
its vanishing endpoint value, which kills the interval jet ladder.

\subsection{Repository normalization}
The operator appearing in the Rvachev RMS analysis is
\begin{equation}
 (\Lop f)(x)=\frac12\left[
 \sin^2\!\left(\frac{\pi x}{2}\right)f(x/2)
 +\cos^2\!\left(\frac{\pi x}{2}\right)f((x+1)/2)
 \right].
 \label{eq:L-definition}
\end{equation}
It is convenient to remove the outer factor and write
\begin{equation}
 \Pop=2\Lop.
 \label{eq:P=2L}
\end{equation}
Thus $\Pop\one=\one$ and $\Pop$ is a positive Markov operator in the
supremum norm.  All structural proofs are first given for $\Pop$ and then
scaled back to $\Lop$.

We regard $C^r(\T)$ as a complex Banach space with norm
\begin{equation}
 \norm f_{C^r}=\sum_{j=0}^{r}\norm{f^{(j)}}_\infty,
 \qquad r\in\N\cup\{0\}.
 \label{eq:Cr-norm}
\end{equation}
Equivalent standard $C^r$ norms give the same spectrum.  On the interval we
use the analogous norm on $C^r([0,1])$.  Essential spectrum always means the
Fredholm essential spectrum,
\begin{equation}
 \spece(T)=\{\lambda\in\C:T-\lambda\id\text{ is not Fredholm}\}.
 \label{eq:fredholm-essential}
\end{equation}

\subsection{Main theorem in the Rvachev case}
\begin{theorem}[Complete $C^r$ spectrum of the Rvachev operator]
\label{thm:main-rvachev}
Let $r\in\N\cup\{0\}$ and let $\Lop$ be \eqref{eq:L-definition}.
On both $C^r(\T)$ and $C^r([0,1])$,
\begin{align}
 \spece(\Lop)&=\Dbar{2^{-r-1}},
 \label{eq:main-essential}\\
 \spec(\Lop)&=\Dbar{2^{-r-1}}\cup\{1/2,-1/4\}.
 \label{eq:main-spectrum}
\end{align}
Every $\lambda$ with $|\lambda|<2^{-r-1}$ is an eigenvalue of infinite
geometric multiplicity.

The following exterior multiplicity statements hold.
\begin{enumerate}[label=\textup{(\alph*)}]
\item If $r\ge1$, then $1/2$ is isolated and simple, with eigenspace
      $\Span\{\one\}$.
\item If $r\ge2$, then $-1/4$ is isolated and semisimple with algebraic and
      geometric multiplicity two.  Its eigenspace is
      \begin{equation}
       \Span\left\{\sin(2\pi x),\ \cos(2\pi x)+\frac13\right\}.
       \label{eq:minus-quarter-space}
      \end{equation}
\end{enumerate}
At $r=1$, the value $-1/4$ lies on the essential boundary; at $r=0$ it lies
strictly inside the essential disk.  The displayed smooth eigenfunctions
remain valid there, but no finite isolated multiplicity is asserted.
\end{theorem}

The threshold picture is therefore exact.
\begin{center}
\small
\begin{tabularx}{0.92\textwidth}{@{}cL{0.18\textwidth}YY@{}}
\toprule
$r$ & Essential radius & Status of $1/2$ & Status of $-1/4$\\
\midrule
$0$ & $1/2$ & essential boundary & interior of essential disk\\
$1$ & $1/4$ & isolated, simple & essential boundary\\
$r\ge2$ & $2^{-r-1}$ & isolated, simple & isolated, semisimple, rank $2$\\
\bottomrule
\end{tabularx}
\end{center}

The finite analytic spectrum from the predecessor is not contradicted.  It
lives on a much smaller function space.  The present disk quantifies exactly
what finite smoothness restores.

\subsection{Sharp operator asymptotics}
The spectral set immediately gives an operator-level consequence that is
stronger than a mere upper bound.

\begin{corollary}[Two-resonance expansion and optimal residual rate]
\label{cor:asymptotic}
On $C^r(\T)$ or $C^r([0,1])$ with $r\ge2$, let $\Pi_+$ and $\Pi_-$ be the
Riesz projections of $\Lop$ at $1/2$ and $-1/4$.  Then
\begin{equation}
 \Lop^n=2^{-n}\Pi_+ +(-1/4)^n\Pi_-+R_n,
 \label{eq:two-resonance-expansion}
\end{equation}
where $\operatorname{rank}\Pi_+=1$, $\operatorname{rank}\Pi_-=2$, and for
every $\varepsilon>0$ there is $C_{r,\varepsilon}$ such that
\begin{equation}
 \norm{R_n}_{C^r\to C^r}
 \le C_{r,\varepsilon}(2^{-r-1}+\varepsilon)^n.
 \label{eq:upper-residual}
\end{equation}
The exponential threshold cannot be improved by any compact correction:
for every compact $K:C^r\to C^r$ and every $n\ge1$,
\begin{equation}
 \norm{\Lop^n-K}_{C^r\to C^r}\ge 2^{-(r+1)n}.
 \label{eq:compact-lower-bound-L}
\end{equation}
In particular, \eqref{eq:compact-lower-bound-L} applies to the finite-rank
resonance sum in \eqref{eq:two-resonance-expansion}.
\end{corollary}

For $r=1$ there is still a leading rank-one decomposition, but the alternating
value $-1/4$ belongs to the essential boundary.  Thus the scalar alternating
RMS law in the repository and the operator-norm geometry genuinely separate
at the critical differentiability order.

\section{A universal dyadic Markov family}
\label{sec:general-family}

\subsection{Weights and circle compatibility}
Let $\omega\in C^\infty(\R)$ satisfy
\begin{equation}
 0\le\omega(x)\le1,
 \qquad
 \omega(x+1)=1-\omega(x).
 \label{eq:weight-hypotheses}
\end{equation}
Define, initially for a $1$-periodic function $f$,
\begin{equation}
 (\Pop_\omega f)(x)=\omega(x)f(x/2)
 +(1-\omega(x))f((x+1)/2).
 \label{eq:general-P}
\end{equation}
Although the two summands are not separately periodic, their sum is:
using \eqref{eq:weight-hypotheses} and periodicity of $f$ gives
$\Pop_\omega f(x+1)=\Pop_\omega f(x)$.  Thus \eqref{eq:general-P} defines a
bounded operator on every $C^r(\T)$.

The Rvachev mask is
\begin{equation}
 \omega_\mathrm{R}(x)=\sin^2(\pi x/2),
 \qquad 1-\omega_\mathrm{R}(x)=\cos^2(\pi x/2).
 \label{eq:rvachev-weight}
\end{equation}
It satisfies $\omega_\mathrm{R}(0)=0$, a fact that will eliminate the
interval endpoint ladder.

\subsection{The universal theorem}
\begin{theorem}[Universal $C^r$ essential disk]
\label{thm:universal-disk}
Assume \eqref{eq:weight-hypotheses}.  For every integer $r\ge0$,
\begin{equation}
 \spece(\Pop_\omega;C^r(\T))=\Dbar{2^{-r}}.
 \label{eq:universal-essential}
\end{equation}
Every $\lambda$ with $|\lambda|<2^{-r}$ is an eigenvalue of
$\Pop_\omega$ on $C^r(\T)$ with infinite geometric multiplicity.
Furthermore:
\begin{enumerate}[label=\textup{(\roman*)}]
\item the powers of $\Pop_\omega$ are uniformly bounded on each $C^m(\T)$;
\item every spectral value with $|\lambda|>2^{-r}$ is an isolated eigenvalue
      of finite algebraic multiplicity;
\item every generalized eigenvector belonging to such a value lies in
      $C^\infty(\T)$.
\end{enumerate}
\end{theorem}

The exact disk, rather than merely its radius, is the key point.  The proof is
split into four transparent steps: a right inverse, an explicit defect space,
a continuous-space disk, and an exact Calkin conjugacy under differentiation.

\section{Right inverse, defect vectors, and the continuous disk}
\label{sec:right-inverse}

\subsection{The doubling right inverse}
Let $T_2x=2x\pmod1$ and define
\begin{equation}
 (\Uop f)(x)=f(T_2x).
 \label{eq:U-definition}
\end{equation}
Then the partition-of-unity identity in \eqref{eq:general-P} gives an exact
one-sided inverse.

\begin{lemma}[Right inverse]
\label{lem:right-inverse}
On every periodic function space on which the formulas make sense,
\begin{equation}
 \Pop_\omega\Uop=\id.
 \label{eq:PU=I}
\end{equation}
Moreover, for $0\le j\le r$ and $n\ge0$,
\begin{equation}
 (\Uop^n f)^{(j)}(x)=2^{nj}f^{(j)}(T_2^n x).
 \label{eq:U-derivatives}
\end{equation}
\end{lemma}
\begin{proof}
Both inverse branches are mapped by $T_2$ back to $x$ modulo one.  Hence
\[
 (\Pop_\omega\Uop f)(x)
 =\omega(x)f(x)+(1-\omega(x))f(x)=f(x).
\]
The derivative identity follows by repeated application of the chain rule.
\end{proof}

\subsection{An explicit infinite-dimensional kernel}
Continuity and \eqref{eq:weight-hypotheses} imply that $\omega$ takes a value
strictly between zero and one.  Choose a nonempty open interval
$I\Subset(0,1)$ on which both branch weights are positive.  For
$g\in C_c^\infty(I)$, define $E_g\in C^\infty(\T)$ by
\begin{equation}
 E_g(y)=
 \begin{cases}
 (1-\omega(2y))g(2y),&y\in I/2,\\
 -\omega(2y-1)g(2y-1),&y\in (I+1)/2,\\
 0,&\text{otherwise}.
 \end{cases}
 \label{eq:defect-vector}
\end{equation}
The two support intervals are disjoint and compactly contained in the two
half-circles, so the extension by zero is smooth.

\begin{lemma}[Branch-cancellation defect]
\label{lem:defect}
For every $g\in C_c^\infty(I)$,
\begin{equation}
 \Pop_\omega E_g=0.
 \label{eq:defect-kernel}
\end{equation}
The map $g\mapsto E_g$ is injective.  Consequently
$\ker\Pop_\omega\cap C^\infty(\T)$ is infinite-dimensional.
\end{lemma}
\begin{proof}
For $x\in I$, the two branch values are
\[
 E_g(x/2)=(1-\omega(x))g(x),\qquad
 E_g((x+1)/2)=-\omega(x)g(x),
\]
so their weighted sum vanishes.  Outside $I$, both values vanish.  Positivity
of $1-\omega$ on $I$ makes the first branch recover $g$, proving injectivity.
\end{proof}

\subsection{A disk of infinite-multiplicity eigenvalues}
For $e\in\ker\Pop_\omega\cap C^\infty(\T)$ and a complex number satisfying
$|\lambda|2^r<1$, consider
\begin{equation}
 F_{\lambda,e}=\sum_{n=0}^\infty \lambda^n\Uop^n e.
 \label{eq:eigenseries}
\end{equation}
By \eqref{eq:U-derivatives}, the series converges absolutely in $C^r$.

\begin{proposition}[Defect eigenseries]
\label{prop:defect-eigenseries}
If $|\lambda|<2^{-r}$, then
\begin{equation}
 \Pop_\omega F_{\lambda,e}=\lambda F_{\lambda,e}.
 \label{eq:eigenseries-equation}
\end{equation}
For fixed $\lambda$, the map $e\mapsto F_{\lambda,e}$ is injective.
Therefore every point of $D(0,2^{-r})$ is an eigenvalue of infinite geometric
multiplicity on $C^r(\T)$.
\end{proposition}
\begin{proof}
Using $\Pop_\omega e=0$ and \eqref{eq:PU=I},
\[
 \Pop_\omega F_{\lambda,e}
 =\sum_{n=1}^\infty\lambda^n\Uop^{n-1}e
 =\lambda F_{\lambda,e}.
\]
Also $(\id-\lambda\Uop)F_{\lambda,e}=e$, so the map is injective.
\end{proof}

At $r=0$, positivity and the partition of unity give
\begin{equation}
 \norm{\Pop_\omega f}_\infty\le\norm f_\infty.
 \label{eq:C0-contraction}
\end{equation}
The preceding eigenseries now determines the entire continuous-space
spectrum and essential spectrum.

\begin{corollary}[Exact continuous-space disk]
\label{cor:C0-disk}
On $C(\T)$,
\begin{equation}
 \spec(\Pop_\omega)=\spece(\Pop_\omega)=\Dbar1.
 \label{eq:C0-disk}
\end{equation}
\end{corollary}
\begin{proof}
The norm bound puts the spectrum in $\Dbar1$.  Every point of the open unit
disk has infinite-dimensional kernel for $\Pop_\omega-\lambda\id$, so it
belongs to the Fredholm essential spectrum.  The non-Fredholm set is closed,
therefore it contains the closed disk.
\end{proof}

\section{Differentiation and the exact Calkin image}
\label{sec:calkin}

\subsection{The leading derivative identity}
For $m\ge1$, Leibniz' rule gives
\begin{equation}
 D^m\Pop_\omega f=2^{-m}\Pop_\omega(D^m f)+\mathcal R_m f,
 \label{eq:derivative-identity}
\end{equation}
where
\begin{equation}
 \mathcal R_m f=
 \sum_{k=1}^{m}\binom{m}{k}2^{-(m-k)}\omega^{(k)}(x)
 \left[f^{(m-k)}(x/2)-f^{(m-k)}((x+1)/2)\right].
 \label{eq:Remainder-m}
\end{equation}
The displayed expression is periodic, even though its individual branch
terms need not be.

\begin{lemma}[Compact lower-jet remainder]
\label{lem:compact-remainder}
For $m\ge1$, $\mathcal R_m:C^m(\T)\to C(\T)$ is compact and annihilates
constants.
\end{lemma}
\begin{proof}
Every term in \eqref{eq:Remainder-m} uses derivatives of order at most
$m-1$.  Thus $\mathcal R_m$ factors through the compact embedding
$C^m(\T)\hookrightarrow C^{m-1}(\T)$, followed by a bounded branch operator.
It annihilates constants because the two branch values then agree.
\end{proof}

\subsection{Removing constants}
Let
\begin{equation}
 Y=\left\{g\in C(\T):\int_0^1g(x)\dd x=0\right\},
 \qquad
 (\Jop g)(x)=g(x)-\int_0^1g(t)\dd t.
 \label{eq:Y-J}
\end{equation}
The map $D^r$ induces a Banach-space isomorphism
\begin{equation}
 \dot D^r:C^r(\T)/\Span\{\one\}\longrightarrow Y.
 \label{eq:derivative-isomorphism}
\end{equation}
Indeed, a continuous zero-mean function can be integrated periodically $r$
times, with integration constants chosen at each stage, and the periodic
Poincare inequality gives boundedness of the inverse.

Define the compressed continuous operator
\begin{equation}
 Q=\Jop\Pop_\omega|_Y:Y\to Y.
 \label{eq:Q-definition}
\end{equation}
With respect to $C(\T)=\Span\{\one\}\oplus Y$, the continuous operator has
an upper triangular block form
\begin{equation}
 \Pop_\omega=
 \begin{pmatrix}
 1&\varphi\\
 0&Q
 \end{pmatrix}
 \label{eq:block-C0}
\end{equation}
for a bounded functional $\varphi$.  Finite-dimensional extensions do not
change Fredholm essential spectrum, so \cref{cor:C0-disk} implies
\begin{equation}
 \spece(Q)=\Dbar1.
 \label{eq:Q-essential}
\end{equation}

\begin{proposition}[Exact Calkin conjugacy]
\label{prop:calkin-conjugacy}
Let $\dot\Pop_{\omega,r}$ be the operator induced by $\Pop_\omega$ on
$C^r(\T)/\Span\{\one\}$.  Under the isomorphism \eqref{eq:derivative-isomorphism},
\begin{equation}
 \dot D^r\dot\Pop_{\omega,r}(\dot D^r)^{-1}
 =2^{-r}Q+K_r,
 \qquad K_r\text{ compact on }Y.
 \label{eq:calkin-conjugacy}
\end{equation}
\end{proposition}
\begin{proof}
The left-hand side applied to the class of $f$ is represented by
$D^r\Pop_\omega f$.  Apply $\Jop$ to
\eqref{eq:derivative-identity}; this does not change the left side because a
periodic derivative has zero mean.  The leading term becomes
$2^{-r}\Jop\Pop_\omega D^rf=2^{-r}Q(D^rf)$, while
$\Jop\mathcal R_rf$ is compact by \cref{lem:compact-remainder}.  It
descends to the quotient because $\mathcal R_r$ annihilates constants.
\end{proof}

\begin{proof}[Proof of the exact disk in \cref{thm:universal-disk}]
The one-dimensional invariant constant space does not change the Fredholm
essential spectrum.  Compact perturbations and similarity preserve it.
Therefore \eqref{eq:calkin-conjugacy} and \eqref{eq:Q-essential} give
\[
 \spece(\Pop_\omega;C^r(\T))
 =2^{-r}\spece(Q)=\Dbar{2^{-r}}.
\]
The infinite-multiplicity point spectrum in the open disk was proved in
\cref{prop:defect-eigenseries}.
\end{proof}

This argument also yields the sharp compact-approximation obstruction.

\begin{corollary}[Calkin lower bound]
\label{cor:Calkin-lower}
For every $r,n\ge0$ with $n\ge1$ and every compact
$K:C^r(\T)\to C^r(\T)$,
\begin{equation}
 \norm{\Pop_\omega^n-K}_{C^r\to C^r}\ge 2^{-rn}.
 \label{eq:P-compact-lower}
\end{equation}
\end{corollary}
\begin{proof}
The Calkin image of $\Pop_\omega$ has spectrum $\Dbar{2^{-r}}$.
By spectral mapping, the spectrum of its $n$th power is
$\Dbar{2^{-rn}}$.  The quotient norm dominates spectral radius, and the
quotient norm of $\Pop_\omega^n$ is the infimum of
$\norm{\Pop_\omega^n-K}$ over compact $K$.
\end{proof}

\section{Exterior eigenvectors are automatically smooth}
\label{sec:bootstrap}

The Calkin theorem determines all non-Fredholm spectral values.  To classify
the remaining spectrum for a particular mask, one needs to know that exterior
eigenvectors cannot hide at finite regularity.  This follows from a simple
bootstrap once power boundedness is established.

\subsection{Uniform power bounds}
\begin{lemma}[Power boundedness on every integer $C^m$]
\label{lem:power-bounded}
For each integer $m\ge0$ there is a constant $B_m$ such that
\begin{equation}
 \norm{\Pop_\omega^n f}_{C^m}\le B_m\norm f_{C^m}
 \qquad(n\ge0).
 \label{eq:power-bounded}
\end{equation}
Consequently the spectral radius of $\Pop_\omega$ on $C^m(\T)$ is one.
\end{lemma}
\begin{proof}
The case $m=0$ is \eqref{eq:C0-contraction}.  Suppose the assertion is known
through order $m-1$.  Applying \eqref{eq:derivative-identity} to
$\Pop_\omega^{n-1}f$ gives
\[
 \norm{D^m\Pop_\omega^n f}_\infty
 \le2^{-m}\norm{D^m\Pop_\omega^{n-1}f}_\infty
 +C_m\norm{\Pop_\omega^{n-1}f}_{C^{m-1}}.
\]
The induction hypothesis bounds the second term uniformly.  Summing the
geometric recurrence bounds the highest derivative; lower derivatives are
already controlled.  The eigenvalue one, carried by constants, shows that
the spectral radius is exactly one.
\end{proof}

\subsection{Resolvent bootstrap}
\begin{proposition}[Smoothness of exterior generalized eigenspaces]
\label{prop:smooth-bootstrap}
Fix $r\ge1$ and $\lambda\in\C$ with $|\lambda|>2^{-r}$.  If
\begin{equation}
 (\Pop_\omega-\lambda\id)^qf=0
 \label{eq:generalized-equation}
\end{equation}
for some $q\ge1$ and $f\in C^r(\T)$, then $f\in C^\infty(\T)$.
\end{proposition}
\begin{proof}
First take $q=1$.  Put $u=D^rf$.  From
\eqref{eq:derivative-identity},
\begin{equation}
 (\lambda\id-2^{-r}\Pop_\omega)u=\mathcal R_rf.
 \label{eq:bootstrap-equation}
\end{equation}
The right side lies in $C^1$, because it only involves derivatives of $f$ of
order at most $r-1$.  By \cref{lem:power-bounded}, the spectrum of
$\Pop_\omega$ on both $C$ and $C^1$ lies in the closed unit disk.  Since
$|\lambda|>2^{-r}$, the operator on the left of
\eqref{eq:bootstrap-equation} is invertible on both spaces.  Its $C^1$
solution and its $C$ solution agree by uniqueness, so $u\in C^1$ and
$f\in C^{r+1}$.

Repeat at orders $r+1,r+2,\ldots$.  The strict inequality persists because
$2^{-m}\le2^{-r}$ for $m\ge r$.

For a generalized eigenvector, induct on $q$.  Set
$g=(\Pop_\omega-\lambda\id)f$.  By induction, $g$ is smooth.  The same
argument applied to the inhomogeneous equation
$\Pop_\omega f-\lambda f=g$ has the additional smooth term $D^mg$ on the
right and again raises regularity one derivative at a time.
\end{proof}

\subsection{Why exterior spectral values are eigenvalues}
The complement of the disk \eqref{eq:universal-essential} is connected and
contains infinity.  On that component, the Fredholm index of
$\Pop_\omega-\lambda\id$ is zero.  Hence every spectral point outside the
disk is an isolated eigenvalue of finite algebraic multiplicity.  Together
with \cref{prop:smooth-bootstrap}, this completes the general assertions of
\cref{thm:universal-disk}.

\section{Fourier collapse for the sine mask}
\label{sec:fourier-collapse}

We now specialize to \eqref{eq:rvachev-weight}.  The finite Fourier support
of the mask forces every exterior generalized eigenvector into a
three-dimensional band.

\subsection{Exact Fourier rule}
Write $e_k(x)=\e^{2\pi\ii kx}$.  Since
\begin{equation}
 \sin^2(\pi x)=\frac12-\frac14e_1(x)-\frac14e_{-1}(x),
 \label{eq:sine-fourier}
\end{equation}
the standard roots-of-unity filter gives
\begin{equation}
 \Fourier{\Pop f}(k)
 =\Fourier f(2k)-\frac12\Fourier f(2k-1)
 -\frac12\Fourier f(2k+1).
 \label{eq:P-fourier-rule}
\end{equation}
The absolute sum of the three coefficients is two.

After $n$ iterations, every input frequency contributing to output frequency
$k$ has the form
\begin{equation}
 q=2^nk+\delta,
 \qquad |\delta|\le2^n-1,
 \label{eq:descendant-frequency}
\end{equation}
and the total absolute coefficient sum is at most $2^n$.  Therefore, if
$|k|\ge2$, then every descendant satisfies
\begin{equation}
 |q|\ge2^n+1.
 \label{eq:descendant-lower}
\end{equation}

\begin{proposition}[Finite Fourier band for exterior chains]
\label{prop:finite-band}
Let $r\ge1$, $|\lambda|>2^{-r}$, and
$(\Pop-\lambda\id)^qf=0$ with $f\in C^r(\T)$.  Then
\begin{equation}
 f\in F_1:=\Span\{e_{-1},e_0,e_1\}.
 \label{eq:F1}
\end{equation}
\end{proposition}
\begin{proof}
By \cref{prop:smooth-bootstrap}, $f$ is smooth.  First suppose $q=1$ and
write $a_k=\Fourier f(k)$.  From $\Pop^nf=\lambda^nf$,
\eqref{eq:descendant-lower}, and rapid Fourier decay, for every $N$ and
$|k|\ge2$,
\[
 |\lambda|^n|a_k|
 \le 2^n\sup_{|q|\ge2^n+1}|a_q|
 \le C_N2^{-n(N-1)}.
\]
Choose $N$ with $2^{1-N}<|\lambda|$ and let $n\to\infty$; then $a_k=0$.

For a generalized chain, let
$g=(\Pop-\lambda\id)f$.  By induction on the chain length, $g\in F_1$.
The band $F_1$ is invariant.  Therefore the high-frequency coefficients of
$\Pop^nf$ still equal $\lambda^na_k$ plus terms obtained from iterates of
$g$, all of which vanish at $|k|\ge2$.  The same estimate gives $a_k=0$.
\end{proof}

\subsection{The exact three-mode matrix}
Using \eqref{eq:P-fourier-rule},
\begin{align}
 \Pop e_{-1}&=-\frac12e_{-1}-\frac12e_0,\\
 \Pop e_0&=e_0,\\
 \Pop e_1&=-\frac12e_0-\frac12e_1.
 \label{eq:mode-actions}
\end{align}
Thus, in the ordered basis $(e_{-1},e_0,e_1)$,
\begin{equation}
 A=
 \begin{pmatrix}
 -1/2&0&0\\
 -1/2&1&-1/2\\
 0&0&-1/2
 \end{pmatrix},
 \qquad
 \det(z\id-A)=(z-1)(z+1/2)^2.
 \label{eq:finite-matrix}
\end{equation}
The matrix is diagonalizable.  Its $1$-eigenspace is generated by $e_0$,
and its $-1/2$ eigenspace is
\begin{equation}
 \left\{ae_{-1}+\frac{a+c}{3}e_0+ce_1:a,c\in\C\right\}.
 \label{eq:minus-half-complex}
\end{equation}
Taking real and imaginary combinations gives the basis in
\eqref{eq:minus-quarter-space} after scaling from $\Pop$ to $\Lop$.
% ed. (2026-09-30): Lean counterparts of the three-mode eigen-identities, which the article mentions without naming them
\begin{ednote}
The eigen-identities of this block are machine-checked pointwise for $\Lop$
in
\path{Analysis/FabiusFunction/Lean/FabiusFunction/RMSTransferEigenfunctions.lean}:
$\Lop\one=\tfrac12\one$ is \path{Fabius.rms_transfer_const_eigen},
$\Lop\sin(2\pi\,\cdot)=-\tfrac14\sin(2\pi\,\cdot)$ is
\path{Fabius.rms_transfer_sin_eigen}, and
$\Lop\bigl(1+3\cos(2\pi\,\cdot)\bigr)=-\tfrac14\bigl(1+3\cos(2\pi\,\cdot)\bigr)$,
three times the second vector of \eqref{eq:minus-quarter-space}, is
\path{Fabius.rms_transfer_one_add_cos_eigen}. The same file proves the
even-frequency case of \eqref{eq:P-fourier-rule},
$\Lop\cos(4\pi m\,\cdot)=\tfrac12\cos(2\pi m\,\cdot)$ and its sine analogue
(\path{Fabius.rms_transfer_cos_even_mode},
\path{Fabius.rms_transfer_sin_even_mode}). These are the identities to
which Layer~1 of the Lean plan in \cref{sec:verification-overview} and
the provenance table in \cref{app:provenance} allude; no spectral statement of this article is
formalized.
\end{ednote}

\begin{proof}[Proof of \cref{thm:main-rvachev} on the circle]
The essential spectrum is \cref{thm:universal-disk}, scaled by one half.
Outside that disk, every spectral value is an eigenvalue and every generalized
eigenvector lies in $F_1$ by \cref{prop:finite-band}.  The finite matrix
\eqref{eq:finite-matrix} therefore supplies all exterior values and their
Jordan structure.  The values $1$ and $-1/2$ of $\Pop$ become $1/2$ and
$-1/4$ for $\Lop$.  The interior infinite-multiplicity statement scales in
the same way.
\end{proof}

\section{Endpoint jets and interval spaces}
\label{sec:endpoint-jets}

The interval problem contains finite-dimensional information absent on the
circle: endpoint derivatives need not match.  For the whole Markov family,
that information is exactly computable.

\subsection{The endpoint quotient}
For $f\in C^r([0,1])$, define the mismatch functionals
\begin{equation}
 \Delta_j(f)=f^{(j)}(1)-f^{(j)}(0),
 \qquad 0\le j\le r.
 \label{eq:jet-mismatch}
\end{equation}
Their common kernel is the closed periodic-jet subspace
\begin{equation}
 X_r^{\mathrm{per}}=
 \{f\in C^r([0,1]):\Delta_j(f)=0\text{ for }0\le j\le r\},
 \label{eq:periodic-subspace}
\end{equation}
which is naturally isomorphic to $C^r(\T)$ and has codimension $r+1$.
It is invariant under $\Pop_\omega$.

The quotient action follows from Leibniz' rule and the compatibility
$\omega^{(k)}(1)=-\omega^{(k)}(0)$ for $k\ge1$.

\begin{lemma}[Exact endpoint-jet recursion]
\label{lem:jet-recursion}
For $0\le j\le r$,
\begin{equation}
 \Delta_j(\Pop_\omega f)
 =2^{-j}\omega(0)\Delta_j(f)
 +\sum_{k=1}^{j}\binom{j}{k}2^{-(j-k)}
   \omega^{(k)}(0)\Delta_{j-k}(f).
 \label{eq:jet-recursion}
\end{equation}
Consequently the induced $(r+1)$-dimensional quotient matrix is triangular
with diagonal
\begin{equation}
 \omega(0),\ \frac{\omega(0)}2,\ \ldots,\ \frac{\omega(0)}{2^r}.
 \label{eq:jet-diagonal}
\end{equation}
\end{lemma}
\begin{proof}
Differentiate \eqref{eq:general-P} $j$ times.  At derivative order $k=0$,
the midpoint terms agree at $x=0$ and $x=1$, while the endpoint terms leave
$2^{-j}\omega(0)\Delta_j(f)$.  For $k\ge1$, the first-branch derivative at
$x=1$ is the negative of that at $x=0$, and the derivative of the second
weight has the opposite sign.  The two midpoint contributions cancel,
leaving exactly the $k$th summand in \eqref{eq:jet-recursion}.
\end{proof}

\begin{theorem}[Endpoint-jet ladder]
\label{thm:endpoint-ladder}
On $C^r([0,1])$,
\begin{align}
 \spece(\Pop_\omega)&=\Dbar{2^{-r}},
 \label{eq:interval-essential}\\
 \spec(\Pop_\omega;C^r([0,1]))
 &=\spec(\Pop_\omega;C^r(\T))
   \cup\{\omega(0)2^{-j}:0\le j\le r\}.
 \label{eq:interval-spectrum-general}
\end{align}
Within the finite quotient, the endpoint values carry the algebraic
multiplicities and internal collisions determined by the triangular jet
matrix \eqref{eq:jet-recursion}.  A collision with a periodic resonance may
also create a cross-block Jordan chain; the theorem asserts the complete
spectral set, not a universal formula for that coupled Jordan structure.
\end{theorem}
\begin{proof}
Relative to a bounded complement of the invariant finite-codimensional
subspace $X_r^{\mathrm{per}}$, the interval operator is block upper
triangular.  Its restriction $A$ is the circle operator, and its quotient
$C$ is the finite matrix of \cref{lem:jet-recursion}.  If both
$A-\lambda\id$ and $C-\lambda\id$ are invertible, the usual triangular
inverse formula shows that the full operator is invertible.  This proves the
inclusion of the interval spectrum in the displayed union.

For the reverse inclusion, first let $\lambda\in\spec(A)$.  If
$|\lambda|\le2^{-r}$, then $\lambda$ belongs to the interval essential
spectrum because a finite-dimensional extension does not change that set.
If $|\lambda|>2^{-r}$, then $A-\lambda\id$ is Fredholm of index zero; when it
is noninvertible it has a nonzero kernel.  The same vector, embedded in the
periodic subspace, lies in the kernel of the full operator.  Next let
$\lambda\in\spec(C)$.  If $A-\lambda\id$ is already noninvertible, the
preceding argument applies.  If it is invertible, choose
$0\ne y\in\ker(C-\lambda\id)$ and solve
$(A-\lambda\id)x=-By$ for the off-diagonal block $B$; then $(x,y)$ is a
nonzero kernel vector of the full block operator.  Hence every value in both
diagonal spectra survives.  Essential spectrum is unchanged by the finite
quotient, proving \eqref{eq:interval-essential} as well.
\end{proof}

\begin{corollary}[Endpoint collapse for the sine mask]
\label{cor:sine-interval}
For \eqref{eq:rvachev-weight}, $\omega_\mathrm{R}(0)=0$, so the quotient jet
matrix is strictly triangular and nilpotent.  Hence the interval and circle
spectra of $\Pop$ coincide.  After division by two, this completes
\cref{thm:main-rvachev} on $C^r([0,1])$.
\end{corollary}

The formula \eqref{eq:interval-spectrum-general} gives a precise answer to
the repository's question about unmatched endpoint jets: they create a
geometric ladder, and in the Rvachev case the endpoint zero collapses the
entire ladder to the already present spectral value zero.

\section{Optimal resonance expansions}
\label{sec:resonance-expansions}

\subsection{Isolation thresholds}
For $\Pop$, the essential radius on $C^r$ is $2^{-r}$.  Thus the eigenvalue
$1$ is exterior exactly for $r\ge1$, and $-1/2$ is exterior exactly for
$r\ge2$.  Scaling by one half gives the thresholds in
\cref{thm:main-rvachev}.  Since the exterior generalized eigenspaces lie in
$F_1$ and the matrix \eqref{eq:finite-matrix} is semisimple, the Riesz
projection ranks are one and two.

\subsection{Proof of the sharp expansion}
\begin{proof}[Proof of \cref{cor:asymptotic}]
For $r\ge2$, isolate the two exterior spectral points by disjoint small
contours.  The corresponding Riesz projections commute with $\Lop$ and have
the ranks read from \eqref{eq:finite-matrix}.  On the complementary invariant
subspace, the spectrum is exactly $\Dbar{2^{-r-1}}$.  The spectral radius
formula gives \eqref{eq:upper-residual} for every $\varepsilon>0$.
The lower bound is \cref{cor:Calkin-lower}, scaled by $2^{-n}$; subtracting
the finite-rank resonance terms is a compact correction.
\end{proof}

\begin{remark}[No operator-norm superconvergence]
The repository obtains scalar asymptotic refinements from special observables
whose Fourier support or dual distribution has extra cancellation.  The lower
bound \eqref{eq:compact-lower-bound-L} shows that such improvement cannot hold
uniformly over the full $C^r$ unit ball after any finite-rank correction.  A
scalar observable may annihilate difficult directions; an operator norm
cannot.
\end{remark}

\section{Comparison with analytic and Sobolev regimes}
\label{sec:comparison}

The same algebraic resonances appear in three distinct regularity geometries,
but the noncompact component changes dramatically.

\begin{center}
\small
\begin{tabularx}{\textwidth}{@{}L{0.24\textwidth}L{0.25\textwidth}YY@{}}
\toprule
Space & Essential spectral set for $\Lop$ & Complete spectral set & Mechanism\\
\midrule
Analytic Hardy disk spaces & none away from $0$; compact/trace-class regime
& $\{0,1/2,-1/4\}$
& analytic compactness and finite Fourier determinant\\
Periodic $H^{s,\beta}$ &
$\Dbar{2^{-s}\sqrt{1+\sqrt{17}}/4}$
& that disk $\cup\{1/2,-1/4\}$
& square-mask coisometry and weighted Fourier geometry\\
Periodic or interval $C^r$ & $\Dbar{2^{-r-1}}$
& that disk $\cup\{1/2,-1/4\}$
& branch defect, derivative quotient, finite Fourier collapse\\
\bottomrule
\end{tabularx}
\end{center}

The Sobolev radius contains the factor
$\sqrt{1+\sqrt{17}}/4$, which comes from the positive eigenfunction of a
squared-mask transfer operator.  The $C^r$ radius does not.  In a supremum-jet
norm, the leading noncompact action is simply the highest derivative composed
with inverse branches, producing $2^{-r}$ before the repository's outer
factor $1/2$.  The two formulas therefore encode different norm geometries,
not conflicting calculations.

At the integer thresholds the contrast is concrete.  In the Sobolev scale,
the leading and subleading isolation thresholds are approximately
$0.1785$ and $1.1785$.  In the $C^r$ scale they occur at the exact integers
$r=1$ and $r=2$.  In particular, $-1/4$ is on the essential boundary of
$C^1$, while it is strictly isolated on $C^2$.

\section{Verification, reproducibility, and formalization route}
\label{sec:verification-overview}

\subsection{Exact finite checks}
The accompanying standard-library Python program checks the algebraic parts
of the argument without floating-point eigensolvers:
\begin{itemize}
\item the rational three-mode matrix, its characteristic polynomial, and the
      displayed eigenvectors;
\item the descendant-frequency bound \eqref{eq:descendant-lower} over a
      finite stress range;
\item the triangular endpoint-jet matrices through order twelve, including
      their exact diagonal ladders and sine-mask nilpotence;
\item the scaling of the spectral thresholds and compact lower bounds listed
      in the summary table.
\end{itemize}
These checks are supplementary.  No finite computation proves the
infinite-dimensional Calkin theorem, the smoothness bootstrap, or the
Fredholm alternative.

\subsection{A staged Lean plan}
A practical formalization can be divided into layers.

\paragraph{Layer 1: branch algebra.}
Define $\Pop_\omega$, the doubling operator $\Uop$, and prove
$\Pop_\omega\Uop=\id$.  Formalize the explicit defect vector with compactly
supported smooth input, or first prove the construction for a fixed smooth
bump.  The existing Fabius/Rvachev files already contain related transfer
eigen-identities.

\paragraph{Layer 2: finite Fourier algebra.}
Formalize \eqref{eq:P-fourier-rule}, invariance of $F_1$, the exact matrix
\eqref{eq:finite-matrix}, and its eigenspaces.  This layer is elementary and
can be completed without a general spectral library.

\paragraph{Layer 3: endpoint jets.}
Define $\Delta_j$ and prove \eqref{eq:jet-recursion} by induction on the
Leibniz formula.  Finite-dimensional triangular-spectrum consequences then
give the endpoint ladder.

\paragraph{Layer 4: Banach and Fredholm infrastructure.}
Formalize the periodic derivative quotient, compactness of
$C^r\hookrightarrow C^{r-1}$, invariance of the Fredholm essential spectrum
under compact perturbations, and the Calkin spectral mapping theorem.  This
is the largest reusable investment.  Once available, the universal disk
statement becomes short.

\paragraph{Layer 5: exterior regularity and Fourier collapse.}
Formalize power boundedness, the resolvent bootstrap, rapid Fourier decay of
smooth periodic functions, and the descendant-frequency estimate.  This
connects the general essential theorem to the complete special spectrum.

No Lean source is included in the present package, and no repository build
was run.  The plan distinguishes elementary algebraic targets from the
substantial functional-analytic library work.

\section{Further research questions and topics}
\label{sec:questions}

\begin{question}[Noninteger H\"older and Zygmund scales]
Determine the complete spectrum on $C^{r,\alpha}(\T)$ for
$0<\alpha<1$, on little-H\"older spaces, and at the Zygmund endpoints.  The
defect eigenseries predicts an essential radius at least
$2^{-(r+\alpha+1)}$ for $\Lop$, and standard Lasota--Yorke estimates suggest
the matching upper radius.  The difficult part is a complete exclusion or
classification of finite-regularity eigenfunctions in the annulus between
successive integer thresholds.
\end{question}

\begin{question}[Boundary point spectrum at $C^1$]
Classify the eigenspace and approximate eigenspaces of $-1/4$ on $C^1$.
The value is simultaneously a smooth eigenvalue and a point of the essential
boundary.  Is its geometric multiplicity finite or infinite?  What is the
local resolvent growth, and which defect vectors generate boundary
approximate eigenfunctions?
\end{question}

\begin{question}[General trigonometric Markov masks]
For smooth trigonometric $\omega$ satisfying
$\omega(x+1)=1-\omega(x)$, classify all spectral values outside
$\Dbar{2^{-r}}$ by a finite Fourier matrix.  Determine a sharp band size in
terms of the degree and dilation, and characterize masks for which all
exterior eigenvalues are semisimple.
\end{question}

\begin{question}[Endpoint-ladder collisions]
In \eqref{eq:interval-spectrum-general}, analyze what happens when an endpoint
value $\omega(0)2^{-j}$ collides with a periodic resonance.  Give necessary
and sufficient conditions for a Jordan chain crossing the periodic and jet
blocks, and derive perturbation expansions through such a collision.
\end{question}

\begin{question}[Weighted endpoint spaces]
Replace the ordinary $C^r$ norm by weights that penalize or permit specified
endpoint singularities.  Which portions of the jet ladder become essential,
and when does the finite quotient turn into an infinite boundary spectrum?
A satisfactory theorem should include power and logarithmic weights.
\end{question}

\begin{question}[Exact Riesz projectors]
Compute the leading and alternating Riesz projectors of the Rvachev operator
as explicit functionals on $C^r$, rather than contour integrals.  The ranges
are known from \eqref{eq:minus-quarter-space}; the open problem is an exact,
well-conditioned description of the dual coefficient functionals and their
relation to the Thue--Morse and Stern distributions in the repository.
\end{question}

\begin{question}[Sharp resolvent and pseudospectral bounds]
Determine the asymptotic size of
$\norm{(\lambda\id-\Lop)^{-1}}$ as $\lambda$ approaches the essential circle,
especially near $-1/4$ in $C^1$.  The explicit right inverse and defect
series may permit two-sided bounds and reveal whether the pseudospectrum is
approximately circular or strongly distorted by the finite resonances.
\end{question}

\begin{question}[Nonlinear expanding maps]
Replace doubling by a smooth nonlinear expanding degree-two map.  The highest
derivative then carries a variable factor $|T'|^{-r}$.  Find hypotheses under
which the full Fredholm essential spectrum, not merely its radius, is an
explicit annulus or disk determined by a weighted composition cocycle.
\end{question}

\begin{question}[Higher-dimensional toral maps]
For an expanding integer matrix on a torus, construct the analogue of the
branch-cancellation defect and determine the $C^r$ essential spectral body.
On a cube, endpoint jets should be replaced by a hierarchy of face, edge, and
corner jets.  Identify the resulting finite spectral ladders and their
collision rules.
\end{question}

\begin{question}[Optimal scalar asymptotics at the operator boundary]
The repository's Rvachev moments exhibit a scalar alternating correction even
when the corresponding eigenvalue lies on an essential boundary.  Classify
the continuous linear functionals on $C^1$ that annihilate enough defect
directions to admit an $o(4^{-n})$ remainder.  Determine whether the known
Fourier cutoff is necessary or merely one sufficient mechanism.
\end{question}

\begin{question}[Proof-producing spectral certificates]
Build a verified interface that accepts a rational trigonometric Markov mask,
constructs its finite Fourier block and endpoint-jet matrix, and returns a
certificate for all exterior resonances once the universal Fredholm theorem
has been formalized.  The output should distinguish isolated eigenvalues,
boundary eigenvalues, and collisions with the endpoint ladder.
\end{question}

\section{Conclusion}
The finite-smoothness question in the Rvachev spectral branch has an exact
integer-order answer.  On $C^r$, analytic spectral collapse is replaced by a
closed disk of radius $2^{-r-1}$, together with the same two algebraic
resonances $1/2$ and $-1/4$.  The disk is not an artifact of an estimate: its
interior consists of infinite-multiplicity eigenvalues, its boundary is forced
by closedness of the Fredholm essential spectrum, and its radius gives a
sharp lower bound against every compact approximation.

The proof reveals a reusable structure.  A dyadic Markov partition gives a
right inverse; two inverse branches give an infinite defect space; and the
highest derivative conjugates the Calkin image by the exact scale $2^{-r}$.
For interval spaces, every unmatched endpoint jet is visible in an explicit
triangular quotient, producing a geometric ladder.  The Rvachev sine mask
vanishes at the endpoint, so that ladder collapses and the circle and interval
spectra agree.

This synthesis links four levels of information that had previously appeared
separately in the repository: branch identities, finite Fourier algebra,
finite-smoothness Fredholm geometry, and resonance asymptotics.  It also
isolates the next genuinely difficult frontier: noninteger H\"older and
endpoint-weighted spaces, where the essential radius is plausible but the
complete exterior and boundary spectral structure is not yet settled here.

\appendix
\section{Fredholm facts used in the proof}
\label{app:fredholm}
For completeness, the functional-analytic inputs are listed precisely.

\begin{enumerate}
\item The set of Fredholm operators between Banach spaces is open, and the
      Fredholm index is locally constant.
\item Compact perturbations preserve Fredholmness and index; equivalently,
      the Fredholm essential spectrum is the spectrum of the Calkin image.
\item Similar operators have the same spectrum and essential spectrum.
\item A finite-dimensional invariant extension does not alter Fredholm
      essential spectrum.  A block triangular operator is invertible whenever
      both diagonal blocks are invertible.  Equality with the union of the
      diagonal spectra requires an additional argument; the finite-quotient
      and index-zero conditions needed here are supplied in the proof of
      \cref{thm:endpoint-ladder}.
\item The spectrum of an element in a unital Banach algebra obeys spectral
      mapping, and its norm dominates its spectral radius.
\item Outside the Fredholm essential spectrum, on a connected component
      containing infinity, a noninvertible operator of index zero has a
      nonzero finite-dimensional kernel.  Isolated spectral values have
      finite-rank Riesz projections.
\item The embedding $C^m(\T)\hookrightarrow C^{m-1}(\T)$ is compact for
      $m\ge1$ by Arzela--Ascoli.
\end{enumerate}
All operator-specific identities needed to invoke these facts are proved in
the body of the article.  Standard references include Kato
\cite{kato}, Nussbaum \cite{nussbaum}, and Hennion \cite{hennion}.

\section{Verification program and recorded scope}
\label{app:verification}
The file \texttt{verify.py} uses only the Python standard library.  Its
recorded run writes \texttt{verification.json}.  The default finite ranges
are deliberately modest and exact.  The program checks:
\begin{center}
\small
\begin{tabular}{lr}
\toprule
Check family & Recorded count or range\\
\midrule
Three-mode matrix identities & all $9$ entries\\
Characteristic polynomial coefficients & all coefficients\\
Displayed eigenvectors & $3$ exact vectors\\
Frequency descendants & levels $0$ through $12$\\
Jet matrices & orders $0$ through $12$\\
Nilpotence for the sine endpoint & orders $0$ through $12$\\
Threshold table & orders $0$ through $12$\\
\bottomrule
\end{tabular}
\end{center}
The finite program does not sample the infinite-dimensional spectrum.  Its
role is to catch transcription, sign, indexing, and scaling errors in the
parts of the proof that reduce to exact finite algebra.

\section{Repository provenance and claim boundaries}
\label{app:provenance}
The comparison target is the public repository
\begin{center}
\url{https://github.com/VladimirReshetnikov/ProveIt}
\end{center}
at inspected main-branch commit
\begin{center}
\texttt{ffddaa8b9c89e7bf027e1442cc6216bb010906d0}.
\end{center}
The relevant current paths and source identities were:
\begin{center}
\small
\begin{tabularx}{\textwidth}{@{}L{0.42\textwidth}Y@{}}
\toprule
Repository item & Role in this article\\
\midrule
\texttt{.../rvachev\_up\_fourier\_decay/}\newline
\texttt{README.md}\newline
{\scriptsize\texttt{6f31c8b9c60d8a2fc10e}\newline
\texttt{b8fb55e8a752672efe42}}
& Group-level status, analytic and Sobolev arrivals, and explicit note that
H\"older spaces remain untreated.\\
\texttt{.../Spectral\_Collapse\_}\newline
\texttt{Alternating\_RMS/article.tex}\newline
{\scriptsize\texttt{cad8f114f2298bd46907c}\newline
\texttt{451f84ace56f93fa784}}
& Source of the question ``Spectra at finite smoothness'' and the analytic
three-point spectrum.\\
\texttt{.../Sobolev\_Spectral\_Disks\_}\newline
\texttt{Rvachev\_Thue\_Morse/article.tex}\newline
{\scriptsize\texttt{ff81881e993a5400489715}\newline
\texttt{8c8e6a11da5ebc1b69}}
& Exact Sobolev disk and explicit statement that H\"older/$C^r$ remains open.\\
\texttt{Analysis/FabiusFunction/Lean/}\newline
\texttt{FabiusFunction/}
& Existing machine-checked transfer eigen-identities and related Fabius
infrastructure; no new build or theorem is claimed here.\\
\bottomrule
\end{tabularx}
\end{center}
The ellipses abbreviate the common path
\begin{center}\small
\path{Analysis/FabiusFunction/docs/semi-formalized-research-frontiers/drafts/}.
\end{center}
The repository was changing rapidly on 30 September 2026.  The immutable
commit and blob identifiers prevent later editorial movement from being
mistaken for the exact comparison corpus.

The literature check located directly relevant work on transfer operators on
the circle with trigonometric sine and cosine weights \cite{chen-volkmer},
general weighted-isometry disk phenomena \cite{bardadyn-kwasniewski}, and the
classical smooth/H\"older transfer-operator framework
\cite{ruelle,hennion}.  This is not an exhaustive worldwide priority search.
Accordingly, the manuscript claims a rigorous resolution of the stated
repository gap and a self-contained structural theorem, not certified global
first publication.

\begingroup
\scriptsize
\raggedright
\raggedcolumns
\setlength{\parskip}{0pt}
\setlength{\columnsep}{1.25em}
\begin{multicols}{2}
\begin{thebibliography}{99}
\setlength{\itemsep}{0pt}

\bibitem{repo-spectral}
Vladimir Reshetnikov,
\emph{Spectral Collapse and Alternating RMS Asymptotics for the Rvachev
Up-Function}, in the \repo{} repository, research manuscript dated
28 September 2026 and editorially updated 29 September 2026.
Inspected at commit
\texttt{ffddaa8b9c89e7bf027e1442cc6216bb010906d0}.

\bibitem{repo-sobolev}
OpenAI ChatGPT,
\emph{From Analytic Spectral Collapse to Sobolev Spectral Disks: Exact
regularity thresholds and sharp endpoint laws for Rvachev--Thue--Morse
transfer operators}, research manuscript prepared for Vladimir Reshetnikov,
29 September 2026, in the \repo{} repository.

\bibitem{chen-volkmer}
Xianghong Chen and Hans Volkmer,
\emph{On transfer operators on the circle with trigonometric weights},
2016, arXiv:1610.07658.

\bibitem{bardadyn-kwasniewski}
Krzysztof Bardadyn and Bartosz Kwa\'sniewski,
\emph{Spectrum of weighted isometries: $C^*$-algebras, transfer operators
and topological pressure}, 2019, revised 2020, arXiv:1911.04811.

\bibitem{ruelle}
David Ruelle,
\emph{The thermodynamic formalism for expanding maps},
Communications in Mathematical Physics \textbf{125} (1989), 239--262.
DOI: 10.1007/BF01217908.

\bibitem{hennion}
Hubert Hennion,
\emph{Sur un th\'eor\`eme spectral et son application aux noyaux
lipchitziens},
Proceedings of the American Mathematical Society \textbf{118} (1993),
627--634.

\bibitem{nussbaum}
Roger D. Nussbaum,
\emph{The radius of the essential spectrum},
Duke Mathematical Journal \textbf{37} (1970), 473--478.

\bibitem{kato}
Tosio Kato,
\emph{Perturbation Theory for Linear Operators},
Springer, corrected printing of the second edition, 1980.

\end{thebibliography}
\end{multicols}
\endgroup

\end{document}
